\subsection{Characters and conjugacy classes}

Let \(\varrho\) be a unitary representation of G in a finite-dimensional Hilbert space E. The character of \(\varrho\) satisfies \(|\chi_{\varrho}| \leqslant \dim(E)\) . Let \((e_i)_{i \in I}\) be an orthonormal basis of E; the character \(\chi_{\varrho}\) is the sum of the diagonal matrix coefficients \(g \mapsto \langle e_i | \varrho(g)e_i \rangle\) of \(\varrho\) . In particular, the character of \(\varrho\) is a G-finite function, and if \(\varrho\) is irreducible, then \(\chi_{\varrho} \in \varrho_{\mathrm{G}}(\varrho)\) .

We have the following formulas

\[
\chi_ {\varrho_ {1} \oplus \varrho_ {2}} = \chi_ {\varrho_ {1}} + \chi_ {\varrho_ {2}}, \quad \chi_ {\varrho_ {1} \otimes \varrho_ {2}} = \chi_ {\varrho_ {1}} \chi_ {\varrho_ {2}}, \quad \chi_ {\breve {\varrho}} = \chi_ {\overline {{\varrho}}} = \overline {{\chi}} _ {\varrho}
\]

(cf. A, VIII, p. 388--389).

PROPOSITION 8. --- We have

\begin{enumerate}
\def\labelenumi{(\arabic{enumi})}
\tightlist
\item
  \(\chi_{\pi}*\chi_{\sigma} = 0\) for all \(\pi, \sigma\) belonging to \(\widehat{\mathrm{G}}\) , \(\pi \neq \sigma\) ,
\end{enumerate}

\[
\chi_ {\pi} * \chi_ {\pi} = \frac {1}{\dim (\pi)} \chi_ {\pi} \quad\text{for every }\pi\in\widehat {\mathrm{G}}.\tag{2}
\]

Let \(\pi\) and \(\sigma\) be irreducible representations of G. We have

\[
\dim (\pi) \left(\chi_ {\pi} * \chi_ {\sigma}\right) = \dim (\pi) \boldsymbol {\gamma} _ {G} \left(\chi_ {\pi}\right) \chi_ {\sigma}
\]

(lemma 4 of V, p. 407). The function \(\dim(\pi)\gamma_{\mathrm{G}}(\chi_{\pi})\chi_{\sigma}\) is the orthogonal projection of \(\chi_{\sigma}\) onto the \(\overline{\pi}\)-isotypic component of \(\gamma_{\mathrm{G}}\) (cor. 2 of V, p. 465), that is, onto \(\varrho_{\mathrm{G}}(\pi)\) (cor. 4 of V, p. 463). Since \(\chi_{\sigma}\) belongs to \(\varrho_{\mathrm{G}}(\sigma)\) , the result follows from theorem 1 of V, p. 462.

Lemma 2. --- The graph of the equivalence relation « \(x \in G\) and \(y \in G\) and \(x\) is conjugate to \(y\) » in G is closed.

Indeed, this graph is the image of the continuous map from the compact space G × G into itself defined by \((x, y) \mapsto (x, yxy^{-1})\) .

We denote by G \(^{\sharp}\) the space of conjugacy classes of G equipped with the quotient topology; it is a compact space by lemma 2 and TG, I, p. 78, prop. 8. Let \(\varpi\) : G → G \(^{\sharp}\) be the canonical projection. The map from \(\mathcal{C}(G^{\sharp})\) into \(\mathcal{C}(G)\) defined by \(f \mapsto f \circ \varpi\) identifies the stellar algebra \(\mathcal{C}(G^{\sharp})\) with the stellar subalgebra of \(\mathcal{C}(G)\) formed by continuous central functions.

The measures on G \(^{\sharp}\) are identified with the central measures on G (V, p. 402, def. 1).

The linear form on \(\mathcal{C}(\mathrm{G}^{\sharp})\) defined by \(f\mapsto\int_{\mathrm{G}}f\) is then a positive measure of mass 1 on \(\mathrm{G}^{\sharp}\) , which is denoted \(\mu^{\sharp}\) . For every \(p\in[1,+\infty]\) , we denote by \(\mathcal{L}^{p}(\mathrm{G}^{\sharp})\) (resp. \(\mathrm{L}^{p}(\mathrm{G}^{\sharp})\) ) the space \(\mathcal{L}_{\mathbf{C}}^{p}(\mathrm{G}^{\sharp},\mu^{\sharp})\) (resp. \(\mathrm{L}_{\mathbf{C}}^{p}(\mathrm{G}^{\sharp},\mu^{\sharp})\) ). We identify \(\mathrm{L}^{p}(\mathrm{G}^{\sharp})\) with the closure of \(\mathcal{C}(\mathrm{G}^{\sharp})\) in \(\mathrm{L}^{p}(\mathrm{G})\) ; in particular, it is a closed subspace of \(\mathrm{L}^{p}(\mathrm{G})\) .

We also denote by \(\Theta(G^{\sharp}) = \Theta(G) \cap \mathcal{C}(G^{\sharp})\) the space of central matrix coefficients of G. It is a unital involutive subalgebra of \(\mathcal{C}(G^{\sharp})\) .

When G is a Lie group, the space \(\Theta(G^{\sharp})\) is also denoted \(Z\Theta(G)\) in LIE, IX, p. 71.

Lemma 3. --- Let \(\pi\) be an irreducible unitary representation of G. The vector space \(\mathrm{L}^2 (\mathrm{G}^\sharp)\cap \varrho_{\mathrm{G}}(\pi)\) is one-dimensional and generated by the character of \(\pi\) .

Consider the unitary representation \(\sigma\) of G in the space \(\varrho_{\mathrm{G}}(\pi)\) defined by \(\sigma(g) = \varrho_{\mathrm{G}}(g, g)\) . Since the unitary representation \(\varrho_{\mathrm{G}}(\pi)\) of G × G is isomorphic to \(\overline{\pi} \boxtimes \pi\) (prop. 5 of V, p. 462), its character is given by

\[
\chi_ {\sigma} (g) = \chi_ {\overline {{\pi}} \boxtimes \pi} (g, g) = | \chi_ {\pi} (g) | ^ {2}
\]

for all \(g \in G\) . The space \(\mathrm{L}^2 (\mathrm{G}^\sharp) \cap \varrho_{\mathrm{G}}(\pi)\) is the subspace of invariant elements of this representation, and its dimension is equal to

\[
\int_ {\mathrm{G}} \chi_ {\sigma} (g) d \mu (g) = \int_ {\mathrm{G}} | \chi_ {\pi} (g) | ^ {2} d \mu (g) = 1
\]

(cor. 3 of V, p. 466 and cor. 3 of V, p. 459). Since the character of \(\pi\) is a nonzero element of \(L^2(G^\sharp) \cap \varrho_G(\pi)\) , it is a basis of this space.

PROPOSITION 9. --- The family \((\chi_{\pi})_{\pi \in \widehat{\mathbf{G}}}\) of characters of irreducible unitary representations of G is an orthonormal basis of \(\mathrm{L}^2 (\mathrm{G}^\sharp)\) .

By the Peter--Weyl theorem (th. 1 of V, p. 462), the closed subspace \(L^{2}(G^{\sharp})\) of \(L^{2}(G)\) formed by the invariant elements of the unitary representation \(g \mapsto \varrho_{\mathrm{G}}(g, g)\) of G is the Hilbert sum of the subspaces \(L^{2}(G^{\sharp}) \cap \varrho_{\mathrm{G}}(\pi)\) for \(\pi \in \widehat{G}\) . The proposition therefore follows from the preceding lemma and cor. 3 of V, p. 459.

COROLLARY 1. --- The vector space \(\Theta(\mathrm{G}^{\sharp})\) is dense in \(\mathcal{C}(\mathrm{G}^{\sharp})\) and in \(\mathrm{L}^{2}(\mathrm{G}^{\sharp})\) .

The second assertion follows from prop. 9. As for the first assertion, consider the linear representation \(\varrho\) of G on the Banach space \(\mathcal{C}(\mathrm{G})\) defined by

\[
\varrho (g) f (x) = f (g ^ {- 1} x g)
\]

for all \(f \in \mathcal{C}(\mathrm{G})\) and every \(x \in \mathrm{G}\) . It is a continuous isometric representation and \(\mathcal{C}(\mathrm{G}^{\sharp})\) is the space of invariant elements of this representation. The continuous projector \(p = \varrho(1)\) of \(\mathcal{C}(\mathrm{G})\) therefore has as its image \(\mathcal{C}(\mathrm{G}^{\sharp})\) (cor. 3 of V, p. 466); it has norm \(\leqslant 1\) .

Let \(f \in \mathcal{C}(\mathrm{G}^{\sharp})\) and let \(\varepsilon > 0\) . There exists \(\widetilde{f} \in \Theta(\mathrm{G})\) such that \(\|f - \widetilde{f}\| \leqslant \varepsilon\) (prop. 4 of V, p. 462), and we then have \(\|f - p(\widetilde{f})\| = \|p(f - \widetilde{f})\| \leqslant \varepsilon\) ; since \(p(\widetilde{f}) \in \mathcal{C}(\mathrm{G}^{\sharp})\) , we conclude that \(\Theta(\mathrm{G}^{\sharp})\) is dense in \(\mathcal{C}(\mathrm{G}^{\sharp})\) .

We denote by R(G) the Grothendieck ring of continuous representations of G in separated finite-dimensional complex vector spaces (cf. LIE, IX, p. 70 and A, VIII, p. 182, applied to the additive class of continuous representations of G in separated finite-dimensional complex vector spaces, viewed as \(\mathbf{C}[G]\)-modules). Since every continuous linear representation of G in a finite-dimensional separated topological vector space is semisimple (prop. 1 of V, p. 457), the abelian group R(G) is free, and the classes of irreducible unitary representations \(\pi \in \widehat{\mathbf{G}}\) form a basis (A, VIII, p. 186, prop. 7).

COROLLARY 2. --- a) The family of characters of the irreducible representations of G is a basis of \(\Theta(G^{\sharp})\) ;

\begin{enumerate}
\def\labelenumi{\alph{enumi})}
\setcounter{enumi}{1}
\tightlist
\item
  The map \(u: \mathrm{R}(\mathrm{G}) \otimes_{\mathbf{Z}} \mathbf{C} \to \Theta(\mathrm{G}^{\sharp})\) such that \(u(\pi \otimes 1) = \chi_{\pi}\) for every \(\pi \in \widehat{\mathrm{G}}\) is an isomorphism of algebras over \(\mathbf{C}\) .
\end{enumerate}

The first assertion follows from Proposition 9 and Corollary 1.

Since the classes of the representations \(\pi \in \widehat{\mathbf{G}}\) form a basis of the free \(\mathbf{Z}\)-module R(G), the map \(u\) is well defined. It is a morphism of \(\mathbf{C}\) -algebras, and it is an isomorphism by \(a\) ).

COROLLARY 3. --- Let H be a locally compact topological group such that G is a compact subgroup of H. Let \(\varrho\) be a unitary representation of H in a Hilbert space E. If the \(\pi\)-isotypic component of the restriction of \(\varrho\) to G is finite-dimensional for all \(\pi \in \widehat{\mathbf{G}}\) , then the representation \(\varrho\) of H is semisimple and every irreducible unitary representation of H has finite multiplicity in \(\varrho\) .

Denote by \(\sigma\) the restriction of \(\varrho\) to the compact subgroup G of H. Let \(\pi\) be an irreducible representation of G. The endomorphism \(\sigma(\chi_{\pi}) \in \mathcal{L}(\mathrm{E})\) has finite rank, since its image is the \(\overline{\pi}\)-isotypic component of \(\sigma\) (cor. 2 of V, p. 465) and it is of finite dimension by hypothesis. Consequently, the endomorphism \(\sigma(f)\) has finite rank for every \(f \in \Theta(\mathrm{G}^{\sharp})\) and is compact for every \(f \in \mathcal{C}(\mathrm{G}^{\sharp})\) (cor. 1 of V, p. 468 and cor. of prop. 2 of III, p. 4).

Denote by \(j\) the canonical injection from G into H. This is a continuous map that is \(\nu\) -proper for every measure \(\nu\) on G (INT, V, p. 69, § 6, no. 1, remark 1). Let \(\mathfrak{B}\) be the filter of compact neighborhoods of the element \(e\) in G. For every V \(\in\) \(\mathfrak{B}\) there exists a continuous positive central function \(f_{\mathrm{V}}\) on G with support contained in V whose integral over G equals 1. Let \(\beta_{\mathrm{V}}\) be the image measure \(j(f_{\mathrm{V}} \cdot \mu)\) ; it is a positive measure with compact support on H such that \(\varrho(\beta_{\mathrm{V}}) = \sigma(f_{\mathrm{V}})\) (INT, V, p. 71, § 6, no. 2, th. 1). According to what precedes, the filter base on \(\mathcal{M}_{c}(\mathrm{H})\) obtained as the image of \(\mathfrak{B}\) under the map \(\mathrm{V} \mapsto \beta_{\mathrm{V}}\) satisfies the conditions of Proposition 2 of V, p. 402, and the assertion follows.

COROLLARY 4 (Weyl's equidistribution criterion)

Let M be a set of central positive measures on G such that \(\nu(G)=1\) for all \(\nu\in M\) . Let \(\mathfrak{F}\) be a filter on M. For the filter \(\mathfrak{F}\) to converge vaguely to the measure \(\mu^{\sharp}\) on \(G^{\sharp}\) , it is necessary and sufficient that, for every nontrivial irreducible unitary representation \(\pi\) , one has

\[
\lim _ {\nu , \mathfrak {F}} \int_ {\mathrm{G}} \chi_ {\pi} (x) d \nu (x) = 0.
\]

Since \(\nu (\mathrm{G}^{\sharp}) = 1 = \mu^{\sharp}(\mathrm{G}^{\sharp})\) for \(\nu \in \mathrm{M}\) , the hypothesis means that

\[
\lim _ {\nu , \mathfrak {F}} \int_ {\mathrm{G} ^ {\sharp}} \chi_ {\pi} (x) d \nu (x) = \int_ {\mathrm{G} ^ {\sharp}} \chi_ {\pi} (x) d \mu^ {\sharp} (x)
\]

for every representation \(\pi\in\widehat{\mathrm{G}}\) , hence it is equivalent to the condition

\[
\lim _ {\nu , \mathfrak {F}} \int_ {\mathrm{G} ^ {\sharp}} f (x) d \nu (x) = \int_ {\mathrm{G} ^ {\sharp}} f (x) d \mu^ {\sharp} (x)
\]

for every function \(f \in \Theta(\mathrm{G}^{\sharp})\) by linearity. Since the space \(\Theta(\mathrm{G}^{\sharp})\) is dense in \(\mathcal{C}(\mathrm{G}^{\sharp})\) (corollary 1), the hypothesis is therefore equivalent to the convergence of the filter \(\mathfrak{F}\) to \(\mu^{\sharp}\) in \(\mathcal{M}(\mathrm{G}^{\sharp})\) endowed with the topology of pointwise convergence in \(\mathcal{C}(\mathrm{G})\) , which coincides with the topology of vague convergence since G is compact (cf. INT, III, p. 59, § 1, no. 9).

\subsection{The Fourier cotransformation}

We equip the set \(\widehat{\mathbf{G}}\) with the discrete topology. We denote by \(\mathrm{F}(\widehat{\mathbf{G}})\) the product algebra of the \(\operatorname{End}(\mathrm{E}_{\pi})\) for \(\pi\) belonging to \(\widehat{\mathbf{G}}\) and \(\mathrm{F}_b(\widehat{\mathbf{G}})\) the product stellar algebra of the \(\operatorname{End}(\mathrm{E}_{\pi})\) (Example 5 of I, p. 103); this is the set of families \((u_{\pi})_{\pi \in \widehat{\mathbf{G}}}\) such that \(\sup \| u_{\pi} \| < +\infty\) .

\[
\pi \in \widehat {\mathbf {G}}
\]

We denote by \(\mathrm{F}_{0}(\widehat{\mathrm{G}})\) the closed stellar subalgebra of \(\mathrm{F}_{b}(\widehat{\mathrm{G}})\) consisting of families \((u_{\pi})_{\pi\in\widehat{\mathrm{G}}}\) such that \(\|u_{\pi}\|\) tends to 0 at infinity.

Let \(\nu \in \mathcal{M}^1 (\mathrm{G})\) . For every \(\pi \in \widehat{\mathrm{G}}\) , we have \(\| \pi (\nu)\| \leqslant \| \nu \|\) , so the family \((\pi (\nu))_{\pi \in \widehat{\mathrm{G}}}\) belongs to \(\mathrm{F}_b(\widehat{\mathrm{G}})\) .

DEFINITION 1. --- For every measure \(\nu \in \mathcal{M}^{1}(\mathrm{G})\) , the element \((\pi (\nu))_{\pi \in \widehat{\mathrm{G}}}\) of \(\mathrm{F}_b(\widehat{\mathrm{G}})\) is denoted \(\overline{\mathcal{F}}_{\mathrm{G}}(\nu)\) . The map from \(\mathcal{M}^1 (\mathrm{G})\) into \(\mathrm{F}_b(\widehat{\mathrm{G}})\) thus defined is called the Fourier cotransformation of G, and \(\overline{\mathcal{F}}_{\mathrm{G}}(\nu)\) is called the Fourier cotransform of the measure \(\nu\) .

For \(f \in \mathrm{L}^1(\mathrm{G})\) , we denote by \(\overline{\mathcal{F}}_{\mathrm{G}}(f) = \overline{\mathcal{F}}_{\mathrm{G}}(f \cdot \mu)\) .

For any representation \(\pi \in \widehat{\mathrm{G}}\) , the map \(\nu \mapsto \pi(\nu)\) is a unital morphism of involutive Banach algebras from \(\mathcal{M}^{1}(\mathrm{G})\) into \(\operatorname{End}(\mathrm{E}_{\pi})\) (lemma 1 of V, p. 401). The Fourier cotransformation is therefore a unital morphism of involutive Banach algebras from \(\mathcal{M}^{1}(\mathrm{G})\) into \(\mathrm{F}_b(\widehat{\mathrm{G}})\) .

\subsection{The Fourier transform}

We keep the notation of the previous number.

Let \(\pi \in \widehat{\mathbf{G}}\) . The vector space \(\operatorname{End}(\mathrm{E}_{\pi})\) is equipped with the structure of a Hilbert space whose scalar product is given by

\[
\langle u _ {1} \mid u _ {2} \rangle = \dim (\pi) \operatorname{Tr} (u _ {1} ^ {*} u _ {2}) = \dim (\pi) \operatorname{Tr} (u _ {2} u _ {1} ^ {*})
\]

for \(u_{1}\) , \(u_{2}\) in \(\operatorname{End}(\mathrm{E}_{\pi})\) (cf. EVT, V, p. 52, th. 1).

We denote by \(\|u\|_{2} = \sqrt{\langle u|u\rangle}\) the norm of an element u of \(\operatorname{End}(\operatorname{E}_{\pi})\) for \(\pi \in \widehat{G}\) . For every \(g \in G\) , we have \(\|\pi(g)\|_{2} = \dim(\pi)\) since \(\pi(g)\) is unitary.

The norm denoted here \(\|u\|_{2}\) differs by a factor \(\dim(\pi)\) from the norm defined in EVT, V, p. 52 on the space of Hilbert–Schmidt operators on \(E_{\pi}\) .

Lemma 4. --- Let \(\pi\) be an irreducible representation of G. The map \(f \mapsto \pi(f)\) defines, by passing to subspaces, an isometric isomorphism from \(\varrho_{\mathrm{G}}(\overline{\pi})\) onto \(\operatorname{End}(\operatorname{E}_{\pi})\) .

Set \(\varepsilon_{\pi}(g,h)u = \pi(g)\circ u\circ\pi(h^{-1})\) for all \((g,h)\in\mathrm{G}\times\mathrm{G}\) and for every \(u\in\mathrm{End}(\mathrm{E}_{\pi})\) . The map \(\varepsilon_{\pi}\) is a continuous representation of \(\mathrm{G}\times\mathrm{G}\) in \(\mathrm{End}(\mathrm{E}_{\pi})\) . It is unitary. Indeed, since \(\pi\) itself is unitary, it follows that

\[
\begin{array}{c} \| \varepsilon_ {\pi} (g, h) u \| _ {2} ^ {2} = \dim (\pi) \operatorname{Tr} \Big ((\pi (g) u \pi (h ^ {- 1})) ^ {*} \pi (g) u \pi (h ^ {- 1}) \Big) \\ = \dim (\pi) \operatorname{Tr} (\pi (h) u ^ {*} u \pi (h) ^ {- 1}) = \dim (\pi) \operatorname{Tr} (u ^ {*} u) = \| u \| _ {2} ^ {2} \end{array}
\]

for all \((g,h)\in \mathrm{G}\times \mathrm{G}\) and every \(u\in \operatorname {End}(\operatorname {E}_{\pi})\) .

The map \(\Psi\) defined by \(f\mapsto\pi(f)\) is a \((G\times G)\) -morphism from \(\varrho_{G}(\overline{\pi})\) into \(\operatorname{End}(E_{\pi})\) since

\[
\pi (\varrho_ {\mathrm{G}} (g, h) f) = \int_ {\mathrm{G}} f (g ^ {- 1} x h) \pi (x) d \mu (x) = \pi (g) \pi (f) \pi (h ^ {- 1})
\]

for all \((g,h)\in\mathrm{G}\times\mathrm{G}\) and every \(f\in\varrho_{\mathrm{G}}(\overline{\pi})\) . Since \(\varrho_{\mathrm{G}}(\overline{\pi})\) is an irreducible representation (th. 1, \(a\) ) of V, p. 462), there exists \(\lambda\in\mathbf{C}^{*}\) such that the map \(\lambda\Psi\) is zero or isometric (cor. 5, \(a\) ) of V, p. 388). Let \(f=\dim(\pi)\overline{\chi}_{\pi}\in\varrho_{\mathrm{G}}(\overline{\pi})\) . We obtain \(\pi(f)=1_{\mathrm{E}_{\pi}}\) (cor. 2 of V, p. 465), whence \(\|\pi(f)\|_{2}=\dim(\pi)=\|f\|\) (cor. 3 of V, p. 459). Consequently, the map \(\Psi\) is an isometry. Since \(\varrho_{\mathrm{G}}(\overline{\pi})\) and \(\operatorname{End}(\mathrm{E}_{\pi})\) have the same dimension, \(\Psi\) is an isometric isomorphism.

We denote by \(\mathrm{F}^{2}(\widehat{\mathrm{G}})\) the Hilbert sum of Hilbert spaces \(\mathrm{End}(\mathrm{E}_{\pi})\) . The norm of an element \(x \in \mathrm{F}^{2}(\widehat{\mathrm{G}})\) is still denoted \(\|x\|_{2}\) .

In LIE, IX, p. 79, this space is denoted \(\mathrm{L}^{2}(\widehat{\mathrm{G}})\) , a notation we prefer to avoid here so as not to create confusion with the space \(\ell^{2}(\widehat{\mathrm{G}})\) . Let \((u_{\pi})_{\pi\in\widehat{\mathrm{G}}}\) be an element of \(\mathrm{F}^{2}(\widehat{\mathrm{G}})\) . Since

\[
\sum_ {\pi \in \widehat {G}} \| u _ {\pi} \| _ {2} ^ {2} = \sum_ {\pi \in \widehat {G}} \dim (\pi) \operatorname{Tr} (u _ {\pi} ^ {*} u _ {\pi}) <   + \infty
\]

and since \(\|u_{\pi}\|^{2}\leqslant\mathrm{Tr}(u_{\pi}^{*}u_{\pi})\) (cf. EVT, V, p. 52, formula (33)), the norm in \(\mathcal{L}(\mathrm{E}_{\pi})\) of the endomorphism \(u_{\pi}\) tends to 0 at infinity. One can therefore identify \(\mathrm{F}^{2}(\widehat{\mathrm{G}})\) with a subspace of \(\mathrm{F}_{0}(\widehat{\mathrm{G}})\) .

Let \(\pi \in \widehat{\mathbf{G}}\) and \(u \in \mathrm{End}(\mathrm{E}_{\pi})\) . We denote by \(\mathcal{F}_{\pi}(u)\) the function on G defined by \(\mathcal{F}_{\pi}(u)(g) = \langle \pi(g) | u \rangle = \dim(\pi) \mathrm{Tr}(\pi(g)^{*} u)\) for all \(g \in \mathbf{G}\) . It is a continuous function on G. If G is a compact real Lie group, then the function \(\mathcal{F}_{\pi}(u)\) is analytic on G (LIE, III, § 8, no. 1, th. 1).

Since G is compact, one may identify \(L^{2}(G)\) with a subspace of \(L^{1}(G)\) .

THEOREM 2. --- The Fourier cotransformation of G induces, by passing to subspaces, an isometric isomorphism of L²(G) onto F²(Ĝ). Its inverse \(\mathcal{F}_{\mathrm{G}}\) associates to an element \((u_{\pi})_{\pi \in \widehat{\mathrm{G}}}\) of F²(Ĝ) the sum of the series

\[
\sum_ {\pi \in \widehat {\mathrm{G}}} \mathcal {F} _ {\pi} (u _ {\pi}),
\]

which converges in \(\mathrm{L}^2 (\mathrm{G})\)

By the Peter–Weyl theorem (Th. 1 of V, p. 462), the Hilbert space \(L^{2}(G)\) is the Hilbert sum of the spaces \(\varrho_{\mathrm{G}}(\overline{\pi})\) for \(\pi \in \widehat{G}\) . For every \(\pi \in \widehat{G}\) , the linear map \(f \mapsto \pi(f)\) from \(\varrho_{\mathrm{G}}(\overline{\pi})\) into \(\operatorname{End}(\operatorname{E}_{\pi})\) is an isometric isomorphism (Lemma 4). Consequently, the restriction to \(L^{2}(G)\) of the Fourier cotransformation defines an isometric isomorphism of \(L^{2}(G)\) onto \(F^{2}(\widehat{G})\) .

Let \(f \in \mathrm{L}^2(\mathrm{G})\) . Let \(\pi \in \widehat{\mathrm{G}}\) and let \(f_{\overline{\pi}} \in \mathcal{C}(\mathrm{G})\) be the orthogonal projection of \(f\) onto \(\varrho_{\mathrm{G}}(\overline{\pi})\) (th. 1 of V, p. 462). Since this space is the \(\overline{\pi}\)-isotypic component of \(\pmb{\delta}_{\mathrm{G}}\) (cor. 4 of V, p. 463), we have \(f_{\overline{\pi}} = \dim(\pi)\pmb{\delta}_{\mathrm{G}}(\chi_{\pi})(f)\) by cor. 2 of V, p. 465. For every \(x \in \mathrm{G}\) , it follows that

\[
\begin{array}{r l} & {f _ {\overline {{\pi}}} (x) = \dim (\pi) \int_ {\mathrm{G}} \chi_ {\pi} (g) (\pmb {\delta} _ {\mathrm{G}} (g) f) (x) d \mu (g)} \\ & {\qquad = \dim (\pi) \int_ {\mathrm{G}} \chi_ {\pi} (g) f (x g) d \mu (g)} \\ & {\qquad = \dim (\pi) \int_ {\mathrm{G}} \chi_ {\pi} (x ^ {- 1} y) f (y) d \mu (y)} \\ & {\qquad = \dim (\pi) \mathrm{Tr} (\pi (x ^ {- 1}) \pi (f)) = \langle \pi (x) | \pi (f) \rangle ,} \end{array}
\]

that is \(f_{\overline{\pi}} = \mathcal{F}_{\pi}(\pi(f))\) .

Since \(f\) is the sum in \(\mathrm{L}^2 (\mathrm{G})\) of the family \((f_{\overline{\pi}})\) , by the Peter–Weyl theorem, we obtain

\[
f = \sum_ {\pi \in \widehat {\mathrm{G}}} \mathcal {F} _ {\pi} (\pi (f))
\]

where the series converges in \(L^{2}(G)\) . This proves the theorem.

DEFINITION 2. --- The isometric isomorphism \(\mathcal{F}_{\mathrm{G}}\) of \(\mathrm{F}^2 (\widehat{\mathrm{G}})\) onto \(\mathrm{L}^2 (\mathrm{G})\) , which is the inverse of the isomorphism induced by the Fourier cotransformation, is called the Fourier transform of G. The image of an element of \(\mathrm{F}^2 (\widehat{\mathrm{G}})\) is called its Fourier transform.

Remarks. --- 1) The Fourier transform of \((u_{\pi})_{\pi \in \widehat{\mathrm{G}}} \in \mathrm{F}^2(\widehat{\mathrm{G}})\) is therefore the class in \(\mathrm{L}^2(\mathrm{G})\) of the series

\[
g \mapsto \sum_ {\pi \in \widehat {\mathrm{G}}} \langle \pi (g) | u _ {\pi} \rangle = \sum_ {\pi \in \widehat {\mathrm{G}}} \dim (\pi) \operatorname{Tr} (u _ {\pi} \circ \pi (g) ^ {- 1}).
\]

\begin{enumerate}
\def\labelenumi{\arabic{enumi})}
\setcounter{enumi}{1}
\tightlist
\item
  Let \(f \in \mathrm{L}^2(\mathrm{G})\) . We then have the Plancherel formula
\end{enumerate}

\[
\| f \| ^ {2} = \| \overline {{\mathcal {F}}} _ {\mathrm{G}} (f) \| ^ {2} = \sum_ {\pi \in \widehat {\mathrm{G}}} \| \pi (f) \| ^ {2}.
\]

Moreover, by Theorem 2, \(f\) is the sum in \(\mathrm{L}^2 (\mathrm{G})\) of the family \((f_{\pi})_{\pi \in \widehat{\mathrm{G}}}\) where

\[
\begin{array}{l} f _ {\pi} (x) = \mathcal {F} _ {\pi} (\pi (f)) (x) \\ \qquad = \langle \pi (x) | \pi (f) \rangle = \dim (\pi) \int_ {\mathrm{G}} f (g) \chi_ {\pi} (x ^ {- 1} g) d \mu (g) \end{array}
\]

for all \(x \in G\) and every \(\pi \in \widehat{G}\) .

Suppose that G is commutative. Since the irreducible representations of G are of dimension 1 (cor. 7 of V, p. 390) and the dual group \(\widehat{\mathbf{G}}\) is discrete (prop. 18 of II, p. 233), the algebras \(\mathrm{F}_b(\widehat{\mathbf{G}})\) and \(\mathrm{F}_0(\widehat{\mathbf{G}})\) are identified, respectively, with the algebra \(\mathcal{C}_b(\widehat{\mathbf{G}})\) of continuous bounded functions on \(\widehat{\mathbf{G}}\) and the algebra \(\mathcal{C}_0(\widehat{\mathbf{G}})\) of continuous functions tending to 0 at infinity on \(\widehat{\mathbf{G}}\) . Since G is compact, the Haar measure on \(\widehat{\mathbf{G}}\) dual to \(\mu\) is the counting measure \(\widehat{\mu}\) (prop. 18 of II, p. 233). The Hilbert sum \(\mathrm{F}^2 (\widehat{\mathrm{G}})\) of the spaces \(\operatorname{End}(\mathrm{E}_{\pi})\) for \(\pi \in \widehat{\mathrm{G}}\) is identified with the Hilbert space \(\mathrm{L}^2 (\widehat{\mathrm{G}},\widehat{\mu})\) .

Let \(\nu \in \mathcal{M}^1 (\mathrm{G})\) . Then, for every unitary character \(\chi \in \widehat{\mathrm{G}}\) , we have

\[
\chi (\nu) = \int_ {\mathrm{G}} \chi \nu
\]

which proves that the Fourier cotransformation defined above coincides with the Fourier cotransformation of G defined in II, p. 206.

Every \(f \in \mathcal{L}^2(\mathrm{G})\) is integrable and the formula \(\| \overline{\mathcal{F}}_{\mathrm{G}}(f)\| _2 = \| f\|\) is the Plancherel formula (II, p. 215, theorem 1).

PROPOSITION 10. --- The Fourier cotransformation is an injective morphism of involutive Banach algebras from \(\mathcal{M}^{1}(\mathrm{G})\) into \(\mathrm{F}_{b}(\widehat{\mathrm{G}})\) which sends \(\mathrm{L}^{1}(\mathrm{G})\) into \(\mathrm{F}_{0}(\widehat{\mathrm{G}})\) .

Since G is compact, the space \(\mathcal{C}(\mathrm{G})\) is contained and dense in \(\mathcal{L}^{1}(\mathrm{G})\) and \(\mathcal{L}^{2}(\mathrm{G})\) .

Let \(\nu \in \mathcal{M}^1 (\mathrm{G})\) such that \(\overline{\mathcal{F}}_{\mathrm{G}}(\nu) = 0\) . For every function \(f\) continuous on G, we then have \(\overline{\mathcal{F}}_{\mathrm{G}}(\nu *f) = 0\) . Now \(\nu *f\) belongs to \(\mathcal{C}(\mathrm{G})\) (INT, VIII, p. 152, § 4, n°2, prop. 3). Since, by theorem 2, the Fourier cotransformation is injective on \(\mathrm{L}^2 (\mathrm{G})\) , and a fortiori, on the space \(\mathcal{C}(\mathrm{G})\) , we have \(\nu *f = 0\) for \(f\in \mathcal{C}(\mathrm{G})\) . In particular, it follows that

\[
\int_ {\mathrm{G}} f (x) d \nu (x) = (\nu * \check {f}) (e) = 0
\]

for every function \(f \in \mathcal{C}(\mathrm{G})\) (INT, VIII, loc. cit.), therefore the measure \(\nu\) is zero.

The image of \(\mathcal{C}(\mathrm{G})\) by the Fourier cotransformation is contained in \(\mathrm{F}^2 (\widehat{\mathrm{G}})\) , and a fortiori in \(\mathrm{F}_0(\widehat{\mathrm{G}})\) . Since \(\mathrm{F}_0(\widehat{\mathrm{G}})\) is closed in \(\mathrm{F}_b(\widehat{\mathrm{G}})\) and since \(\mathcal{C}(\mathrm{G})\) is dense in \(\mathrm{L}^1 (\mathrm{G})\) , the image of \(\mathrm{L}^1 (\mathrm{G})\) by the Fourier cotransformation is also contained in \(\mathrm{F}_0(\widehat{\mathrm{G}})\) .

PROPOSITION 11. --- a) Let \(u = (u_{\pi})_{\pi \in \widehat{\mathbf{G}}}\) be an element of \(\mathrm{F}(\widehat{\mathbf{G}})\) . If the family \((\mathcal{F}_{\pi}(u_{\pi}))_{\pi \in \widehat{\mathbf{G}}}\) is uniformly summable in \(\mathcal{C}(\mathbf{G})\) , then its sum \(f\) is a continuous function on \(\mathbf{G}\) whose Fourier cotransform is \(u\) ;

\begin{enumerate}
\def\labelenumi{\alph{enumi})}
\setcounter{enumi}{1}
\tightlist
\item
  Let \(f \in \mathrm{L}^1(\mathrm{G})\) . If the family \((\mathcal{F}_\pi(\pi(f)))_{\pi \in \widehat{\mathrm{G}}}\) is uniformly summable in \(\mathcal{C}(\mathrm{G})\) , then its sum is a continuous function on \(\mathrm{G}\) whose class in \(\mathrm{L}^1(\mathrm{G})\) is equal to \(f\) .
\end{enumerate}

Let us prove assertion \(a\) ). Since G is compact, the sum \(f\) of the family \((\mathcal{F}_{\pi}(u_{\pi}))\) belongs to \(\mathrm{L}^2 (\mathrm{G})\) and the series

\[
\sum_ {\pi \in \widehat {\mathrm{G}}} \mathcal {F} _ {\pi} (u _ {\pi})
\]

converges in \(L^{2}(G)\) to f (INT, IV, p. 127, § 3, n° 3, prop. 4). We therefore have \(\mathcal{F}_{\mathrm{G}}((u_{\pi})_{\pi}) = f\) in \(L^{2}(G)\) , whence \((u_{\pi})_{\pi \in \widehat{\mathrm{G}}} = \overline{\mathcal{F}}_{\mathrm{G}}(f)\) (theorem 2).

Let us prove assertion b). We may apply assertion a) to the family \((\pi(f))_{\pi \in \widehat{\mathrm{G}}}\) . The sum g of the family \((\mathcal{F}_{\pi}(\pi(f)))_{\pi \in \widehat{\mathrm{G}}}\) is a continuous function, and hence belongs to \(L^{2}(G)\) , such that \(\overline{\mathcal{F}}_{\mathrm{G}}(g) = (\pi(f))_{\pi \in \widehat{\mathrm{G}}}\) . Since \(\overline{\mathcal{F}}_{\mathrm{G}}\) is injective on \(L^{2}(G)\) , we have f = g in \(L^{2}(G)\) , hence in \(L^{1}(G)\) .

The algebra \(\mathcal{C}(\widehat{\mathrm{G}})\) of complex-valued functions on \(\widehat{\mathrm{G}}\) is identified with the center of the algebra \(\mathrm{F}(\widehat{\mathrm{G}})\) by the map that associates to \(f\colon \widehat{\mathrm{G}}\to \mathbf{C}\) the central element \((f(\pi)1_{\mathrm{E}_{\pi}})_{\pi \in \widehat{\mathrm{G}}}\) of \(\mathrm{F}(\widehat{\mathrm{G}})\) . Denote by \(\mathcal{M}^1 (\mathrm{G}^\sharp)\) the space of central bounded measures on G. It is a closed subspace of the Banach algebra \(\mathcal{M}^1 (\mathrm{G})\) . For every measure \(\nu \in \mathcal{M}^1 (\mathrm{G}^\sharp)\) and every representation \(\pi \in \widehat{\mathrm{G}}\) , we have \(\pi (\nu)\in \mathrm{End}_{\mathrm{G}}(\mathrm{E}_{\pi})\) hence \(\pi (\nu)\) is a scalar multiple of the identity endomorphism of \(\mathrm{E}_{\pi}\) by Schur's lemma (prop. 6 of V, p. 386). The restriction of the Fourier cotransformation to \(\mathcal{M}^1 (\mathrm{G}^\sharp)\) is therefore identified with a unital morphism of involutive Banach algebras from \(\mathcal{M}^1 (\mathrm{G}^\sharp)\) into \(\mathcal{C}(\widehat{\mathrm{G}})\) . This morphism is injective; the image of \(\mathrm{L}^1 (\mathrm{G}^\sharp)\) is contained in \(\mathcal{C}_b(\widehat{\mathrm{G}})\) and its restriction to \(\mathrm{L}^2 (\mathrm{G}^\sharp)\) is an isometric isomorphism onto \(\mathrm{L}^2 (\widehat{\mathrm{G}},\widehat{\mu})\) .

\subsection{Frobenius--Schur indicator and Larsen's alternative}

Recall (LIE, VIII, §7, n° 6, def. 2) that an irreducible representation of G in a finite-dimensional complex vector space E is said to be of orthogonal type (resp. of symplectic type) if there exists a nonzero symmetric bilinear form on E invariant under G (resp. if there exists a nonzero alternating bilinear form on E invariant under G). It is said to be of complex type if there exists no nonzero bilinear form on E invariant under G.

When an irreducible representation in E is of orthogonal type (resp. symplectic type), the space of symmetric (resp. alternating) G-invariant bilinear forms on E has dimension 1 and every nonzero G-invariant bilinear form on E is separating.

A representation of orthogonal type is sometimes said to be of real type, and a representation of symplectic type is sometimes said to be of quaternionic type (cf. LIE, IX, App. II).

Let \(\pi\) be an irreducible representation of G in a complex vector space E. The three possibilities above are distinguished by the value of the quantity

\[
\mathrm{FS} (\pi) = \int_ {\mathrm{G}} \chi_ {\pi} (g ^ {2}) d \mu (g),
\]

which is called the Frobenius--Schur indicator of π. Indeed we have

\[
\operatorname{FS} (\pi) = \left\{ \begin{array}{l l} 1 & \text{if }\pi\text{ is of orthogonal type,} \\ - 1 & \text{if }\pi\text{ is of symplectic type,} \\ 0 & \text{if }\pi\text{ is of complex type.} \end{array} \right.
\]

(LIE, IX, App. 2, p. 103, prop. 3 and p. 105, prop. 4).

When G is a Lie group, one can compute \(\mathrm{FS}(\pi)\) using prop. 1 of LIE, IX, p. 69.

DEFINITION 3. --- Let \(\varrho\) be a finite-dimensional unitary representation of G in a Hilbert space E. Let \(k\) be a positive integer. We call the dimension of the subspace of G-invariant elements in the representation \(\overline{\varrho}^{\otimes k} \otimes \varrho^{\otimes k}\) the absolute moment of order \(2k\) of \(\varrho\), and denote it by \(\mathrm{M}_{2k}(\varrho)\) .

THEOREM 3 (Larsen's alternative). --- Suppose that G is infinite. Let π be a faithful unitary representation of G in a Hilbert space E of finite dimension ≥ 2.

\begin{enumerate}
\def\labelenumi{\alph{enumi})}
\item
  Suppose that the derived group of G is infinite. We have \(M_{4}(\pi) \geqslant 2\) with equality if and only if G contains SU(E);
\item
  Suppose that \(\dim(E) \geqslant 3\) , that \(\pi\) is of orthogonal or symplectic type, and that \(\dim(E) \neq 4\) if \(\pi\) is of orthogonal type. Then \(M_4(\pi) \geqslant 3\) with equality if and only if the complexified Lie algebra of G is equal to the Lie algebra of a separating G-invariant bilinear form on E. *This is the case if and only if the neutral component \(G_0\) of G is a maximal compact subgroup of the group of automorphisms of such a bilinear form.*
\end{enumerate}

We shall use in the proof the following lemmas.

Lemma 5. --- Let \(\pi\) be a unitary representation of G in a Hilbert space E of finite dimension. Let \((\varrho_{i})_{i\in\mathrm{I}}\) be a family of nonzero unitary representations of G and \((n_{i})_{i\in\mathrm{I}}\) a family of integers \(\geqslant 1\) such that the representation \(\overline{\pi}\otimes\pi\) of G (resp. the representation \(\pi\otimes\pi\) ) is isomorphic to the direct sum \(\bigoplus_{i\in\mathrm{I}}\varrho_{i}^{n_{i}}\) . Then we have

\[
\mathrm{M} _ {4} (\pi) \geqslant \sum_ {i \in \mathrm{I}} n _ {i} ^ {2}
\]

with equality if and only if the representations \(\varrho_{i}\) are irreducible and pairwise non-isomorphic.

Denote by \(\chi_{i}\) the character of \(\varrho_{i}\) for \(i\in I\) . We have

\[
| \chi_ {\pi} | ^ {2} = \chi_ {\overline {{\pi}} \otimes \pi} = \sum_ {i \in \mathrm{I}} n _ {i} \chi_ {i}, \quad (\text { resp. } \chi_ {\pi} ^ {2} = \chi_ {\pi \otimes \pi} = \sum_ {i \in \mathrm{I}} n _ {i} \chi_ {i}).
\]

By the definition and cor. 3 of V, p. 466, we have

\[
\mathrm{M} _ {4} (\pi) = \int_ {\mathrm{G}} | \chi_ {\pi} | ^ {4} d \mu ,
\]

whence, in both cases, the formula

\[
\mathrm{M} _ {4} (\pi) = \sum_ {i, j} n _ {i} n _ {j} \int_ {\mathrm{G}} \chi_ {i} \overline {{\chi}} _ {j} d \mu = \sum_ {i, j} n _ {i} n _ {j} \dim \mathrm{Hom} _ {\mathrm{G}} (\varrho_ {i}, \varrho_ {j})
\]

(cor. 4, b) of V, p. 459). Let \(i \in I\) . Since the representation \(\varrho_{i}\) is not zero, the space \(\mathrm{Hom}_{\mathrm{G}}(\varrho_{i}, \varrho_{i})\) is not zero, and we deduce the lower bound

\[
\mathrm{M} _ {4} (\pi) \geqslant \sum_ {i \in \mathrm{I}} n _ {i} ^ {2}.
\]

Moreover, since \(n_{i} \geqslant 1\) , there is equality if and only if we have \(\dim \operatorname{Hom}_{\mathrm{G}}(\varrho_{i}, \varrho_{j}) = 0\) if \(i \neq j\) and \(\dim \operatorname{Hom}_{\mathrm{G}}(\varrho_{i}, \varrho_{i}) = 1\) for every i. The second condition holds if and only if the representations \(\varrho_{i}\) are irreducible (loc. cit.). The first is then satisfied if and only if the representations \(\varrho_{i}\) are pairwise non-isomorphic, by Schur's lemma (V, p. 387, cor. 2).

If E is a module over a commutative ring A, we will call the symmetric square (resp. exterior square) of E the A-module \(S^{2}(E)\) (resp. \(\wedge^{2}E\) ).

Lemma 6. --- Let \(q\) be a separating bilinear form on a finite-dimensional complex vector space E of dimension \(\geqslant 3\) .

\begin{enumerate}
\def\labelenumi{\alph{enumi})}
\item
  If \(q\) is symmetric, then the adjoint representation of the orthogonal Lie algebra \(\mathfrak{so}(q)\) is isomorphic to the exterior square \(\wedge^2\mathrm{E}\) of the natural representation \(\mathfrak{so}(q) \to \mathfrak{gl}(\mathrm{E})\) ;
\item
  If q is alternating, then the adjoint representation of the symplectic Lie algebra \(\mathfrak{sp}(q)\) is isomorphic to the symmetric square \(\mathsf{S}^2 (\mathrm{E})\) of the natural representation \(\mathfrak{sp}(q)\to \mathfrak{gl}(\mathrm{E})\) .
\end{enumerate}

Endow C with the identity involutive automorphism. The bilinear form q is then \(\varepsilon\) -Hermitian (A, IX, § 3, no. 1, def. 1) with \(\varepsilon = 1\) in case a) and \(\varepsilon = -1\) in case b). For every \(u \in \text{End}(\text{E})\) , we denote by \({}^{t}u\) the adjoint of u with respect to q. We therefore have \(q(x, u(y)) = q(^{t}u(x), y)\) for all \((x, y) \in \text{E} \times \text{E}\) . Let v denote the unique automorphism of E \(\otimes\) E such that \(v(x \otimes y) = y \otimes x\) for all \((x, y) \in \text{E}^{2}\) .

There exists a unique linear map \(w\) from \(\mathrm{E} \otimes \mathrm{E}\) into \(\operatorname{End}(\mathrm{E})\) such that \(w(a \otimes b)\) is the linear map defined by \(x \mapsto q(a, x)b\) for all \((a, b, x) \in \mathrm{E}^3\) . The map \(w\) is an isomorphism. For every \(s \in \mathrm{E} \otimes \mathrm{E}\) , we have

\[
{ } ^ { t } ( w ( s ) ) = \varepsilon w ( v ( s ) ) .\tag{3}
\]

Indeed, for all \((x,y)\) and \((a,b)\) in \(\mathrm{E} \times \mathrm{E}\) , it follows that

\[
q (x, w (a \otimes b) y) = q (a, y) q (x, b) = \varepsilon q (b, x) q (a, y) = q (\varepsilon w (b \otimes a) x, y).
\]

Finally, denote by \(x \mapsto x^*\) the isomorphism from E to \(\mathrm{E}'\) deduced from \(q\) , that is, such that \(\langle x^*, y \rangle = q(x, y)\) for all \((x, y) \in \mathrm{E}^2\) .

With these notations established, let H denote the orthogonal (resp. symplectic) group of q. It is a Lie subgroup of \(\mathbf{GL}(\mathrm{E})\) whose Lie algebra h is the subspace of End(E) consisting of the \(u \in \operatorname{End}(\operatorname{E})\) such that \(^{t}u = -u\) (LIE, III, p. 146, cor. 1). Since the form q is H-invariant, it follows that

\[
\langle (g a) ^ {*}, b \rangle = q (g a, b) = \langle a, g ^ {- 1} b \rangle
\]

for all \(g \in H\) and every \((a, b) \in E \times E\) .

Equip End(E) with the adjoint representation Ad of H. The linear map \(w\) is a morphism of representations of H: indeed, for every \(g \in \mathrm{H}\) and every \((a, b, x) \in \mathrm{E}^3\) , we have

\[
\begin{array}{c} w (g   (a \otimes b)) (x) = \langle (g a) ^ {*}, x \rangle   g b = \langle a, g ^ {- 1} x \rangle   g b \\ = g   (\langle a, g ^ {- 1} x \rangle   b) = \mathrm{Ad} (g) (w (a \otimes b)) x, \end{array}
\]

whence the conclusion by linearity.

Since v is also a morphism of representations of H, the same holds for the linear map \(\theta = w - \varepsilon w \circ v\) ; this is therefore a morphism of h-modules.

Denote by \(F \subset E \otimes E\) the subspace of elements \(s \in E \otimes E\) such that \(v(s) = -\varepsilon s\) . If \(\varepsilon = -1\) , the restriction to F of the canonical map from \(E \otimes E\) into the symmetric square of E is an isomorphism of C-vector spaces (A, III, p. 72); if \(\varepsilon = 1\) , the restriction to F of the canonical map from \(E \otimes E\) into the exterior square of E is an isomorphism of C-vector spaces (loc. cit., p. 82).

The image of F under \(\theta\) is contained in \(\mathfrak{h}\) : indeed, for every \(s\in\mathrm{F}\) the formula (3) implies

\[
{ } ^ { t } \theta ( s ) = { } ^ { t } w ( s ) - \varepsilon { } ^ { t } w ( v ( s ) ) = \varepsilon w ( v ( s ) ) - w ( s ) = - \theta ( s ) .
\]

By definition of F, the restriction of the map \(\theta\) to F coincides with that of \(2w\) and the linear map \(\theta\) is therefore injective on F. Since \(\dim(F) = \dim(h)\) . From (A, III, p. 75, th. 1 and p. 87, cor. 1 and LIE, VIII, p. 192 and p. 201), we conclude that \(\theta\) induces, by passage to subspaces, an isomorphism of \(h\)-modules from F onto \(h\), whence the lemma.

Lemma 7. --- Let \(H_{2}\) be a real Lie group and let \(H_{1}\) be a closed subgroup of \(H_{2}\) . The complexified Lie algebra of \(H_{1}\) is identified with a subrepresentation of the adjoint representation of \(H_{1}\) on the complexified Lie algebra of \(H_{2}\) .

Indeed, the restriction to \(H_{1}\) of the adjoint representation of \(H_{2}\) on its Lie algebra identifies with the adjoint representation of \(H_{1}\) .

Lemma 8. --- Let E be a complex Hilbert space of finite dimension n.

\begin{enumerate}
\def\labelenumi{\alph{enumi})}
\item
  The groups \(\mathbf{SU}(\mathrm{E})\) and \(\mathbf{U}(\mathrm{E})\) are connected;
\item
  The derived group of SU(E) is equal to SU(E);
\item
  The derived subgroup of U(E) is equal to SU(E).
\end{enumerate}

We may assume that \(n \geqslant 1\) and that \(E = C^{n}\) .

Let A be the subgroup of \(\mathbf{U}(\mathrm{E})\) formed by diagonal matrices; it is homeomorphic to \(U^{n}\) and therefore connected (TG, VI, p. 11, cor. 2 and I, p. 83, prop. 8). Its intersection with \(\mathbf{SU}(\mathrm{E})\) is homeomorphic to \(U^{n-1}\) , hence it is also connected.

By th. 1 of IV, p. 149, the group \(\mathbf{U}(\mathrm{E})\) (resp. \(\mathbf{SU}(\mathrm{E})\) ) is the union of the connected subsets \(g\mathrm{Ag}^{-1}\) for \(g \in \mathrm{G}\) (resp. of the connected subsets \(g(\mathrm{A} \cap \mathbf{SU}(\mathrm{E}))g^{-1}\) ); since the identity element belongs to each of these sets, the space \(\mathbf{U}(\mathrm{E})\) (resp. \(\mathbf{SU}(\mathrm{E})\) ) is connected (TG, I, p. 81, prop. 2).

Let us prove \(b\) ). The assertion is true when \(n = 1\) since \(\mathbf{SU}(\mathrm{E})\) is then reduced to the identity element. Suppose therefore that \(n \geqslant 2\) . The Lie algebra of \(\mathbf{SU}(\mathrm{E})\) is then simple (LIE, IX, p. 20, § 3, n°4), so the derived group of \(\mathbf{SU}(\mathrm{E})\) has finite index in \(\mathbf{SU}(\mathrm{E})\) . From this we deduce the result since \(\mathbf{SU}(\mathrm{E})\) is connected by \(a\) ).

The assertion \(c\) follows from \(b\) since the derived group of \(\mathbf{U}(\mathrm{E})\) is contained in \(\mathbf{SU}(\mathrm{E})\) .

We now prove Theorem 3. Let n denote the dimension of the Hilbert space E. Since the representation \(\pi\) is faithful, one may assume that G is a compact subgroup of U(E). The group G is closed in the real Lie group U(E), hence is a compact real Lie group (LIE, III, §8, no. 2, th. 2). By hypothesis, G is infinite, hence of dimension \(\geqslant 1\) . Let g be its Lie algebra; it is nonzero.

Let us prove \(a\) ). Suppose that the derived group D(G) is infinite. We have the G-invariant decomposition \(\operatorname{End}(\mathrm{E}) = \mathbf{C}1_{\mathrm{E}} \oplus \mathfrak{sl}(\mathrm{E})\) . The representations of G on \(\mathbf{C}1_{\mathrm{E}}\) and on \(\mathfrak{sl}(\mathrm{E})\) are not isomorphic, since \(\dim \mathfrak{sl}(\mathrm{E}) \geqslant 2\) . By Lemma 5, we therefore have \(\mathrm{M}_4(\pi) \geqslant 2\) with equality if and only if the representation of G in \(\mathfrak{sl}(\mathrm{E})\) is irreducible. Now, the derived group of G is contained in \(\mathbf{SL}(\mathrm{E})\) , and the complexification of the Lie algebra of D(G) therefore identifies with a subspace of \(\mathfrak{sl}(\mathrm{E})\) , which is a subrepresentation of the representation of G in \(\mathfrak{sl}(\mathrm{E})\) (Lemma 7). This subrepresentation is nonzero, since D(G) is infinite. We therefore have \(\mathrm{M}_4(\pi) = 2\) if and only if \((\mathcal{D}\mathfrak{g})_{(\mathbf{C})} = \mathfrak{sl}(\mathrm{E})\) . Since D(G) is contained in \(\mathbf{SU}(\mathrm{E})\) , and since \(\mathbf{SU}(\mathrm{E}) = \mathrm{D}(\mathbf{SU}(\mathrm{E}))\) (Lemma 8), this condition is equivalent to \(\mathbf{SU}(\mathrm{E}) \subset \mathbf{G}\) .

For the proof of assertion \(b\) ), we resume the notation of LIE, VIII, §13, nos. 2--4.

Suppose that \(\pi\) is of symplectic type and that \(\dim(\mathrm{E}) \geqslant 3\) . Let \(q\) be a nonzero G-invariant alternating bilinear form on E; it is separating (LIE, IX, p. 103, App. 2, prop. 3) and G is a compact subgroup of the symplectic group \(\mathbf{Sp}(q)\) . Denote by \(q^{*} \in \wedge^{2}\mathrm{E}\) the element to which the alternating form \(q\) identifies. The complexified Lie algebra \(\mathfrak{g}_{(\mathbf{C})}\) of G is contained in \(\mathfrak{sp}(q)\) . Now, since \(\dim(\mathrm{E}) \geqslant 3\) , the representation of \(\mathfrak{sp}(q)\) in \(\mathrm{E} \otimes \mathrm{E}\) admits the direct sum decomposition

\[
\mathrm{E} \otimes \mathrm{E} = \mathsf {S} ^ {2} (\mathrm{E}) \oplus \wedge^ {2} \mathrm{E} = \mathsf {S} ^ {2} (\mathrm{E}) \oplus \mathrm{E} _ {2} \oplus \mathbf {C} q ^ {*},
\]

where \(E_{2}\) is the second fundamental representation of \(\mathfrak{sp}(q)\) (LIE, VIII, p. 202, §13, no. 3, (IV)). The representations \(E_{2}\) and \(C q^{*}\) of \(\mathfrak{sp}(q)\) are irreducible (loc. cit.) and the representation \(S^{2}(E)\) is nonzero.

By Lemma 5, we therefore have \(\mathrm{M}_{4}(\pi) \geqslant 3\) with equality if and only if the representations of G in \(\mathsf{S}^{2}(\mathrm{E})\) , \(E_{2}\) and \(Cq^{*}\) are irreducible and pairwise non-isomorphic.

Suppose that \(M_{4}(\pi)=3\) . Since the representation \(S^{2}(E)\) identifies with the adjoint representation of \(\mathfrak{sp}(q)\) (Lemma 6), it contains as a subrepresentation the complexification of the adjoint representation of G (Lemma 7). We therefore have \(\mathfrak{g}_{(\mathbf{C})}=\mathfrak{sp}(q)\) .

Conversely, suppose that \(\mathfrak{g}_{(\mathbf{C})} = \mathfrak{sp}(q)\) . The adjoint representation of \(\mathfrak{sp}(q)\) and the representation \(\mathrm{E}_2\) of \(\mathfrak{sp}(q)\) are irreducible, of dimensions \(n(n + 1)/2\) and \(n(n - 1)/2 - 1\) respectively (LIE, VIII, p. 202, §13, no. 3, (IV)), which are different and \(\geqslant 2\) since \(n \geqslant 3\) , whence \(\mathrm{M}_4(\pi) = 3\) .

Finally, suppose that \(\pi\) is of real type and that \(\dim(\mathrm{E}) \geqslant 3\) and \(\dim(\mathrm{E}) \neq 4\) . Let \(q\) be a nonzero symmetric bilinear form on \(\mathrm{E}\), invariant under G ; it is separating (LIE, IX, p. 103, App. 2, prop. 3) and \(\mathrm{G}\) is a compact subgroup of the orthogonal group \(\mathbf{O}(q)\) . Denote by \(q^{*} \in \mathbb{S}^{2}(\mathrm{E})\) the element to which \(q\) identifies.

The complexified Lie algebra \(\mathfrak{g}_{(\mathbf{C})}\) of G is contained in \(\mathfrak{so}(q)\) . The representation of \(\mathfrak{so}(q)\) in E \(\otimes\) E admits the direct sum decomposition

\[
\mathrm{E} \otimes \mathrm{E} = \wedge^ {2} \mathrm{E} \oplus \mathsf {S} ^ {2} (\mathrm{E}) = \wedge^ {2} \mathrm{E} \oplus \mathsf {S} _ {0} ^ {2} (\mathrm{E}) \oplus \mathbf {C} q ^ {*},
\]

where \(\mathsf{S}_{0}^{2}(\mathbf{E})\) is the orthogonal complement of \(q^{*}\) in \(\mathsf{S}^{2}(\mathbf{E})\) . These representations are of dimension at least 2 since \(\dim(\mathbf{E})\geqslant3\) . By Lemma 5, we have \(\mathrm{M}_{4}(\pi)\geqslant3\) with equality if and only if the representations of G in \(\wedge^{2}\mathrm{E}\) and \(\mathsf{S}_{0}^{2}(\mathbf{E})\) are irreducible and non-isomorphic.

Suppose that \(\mathrm{M}_{4}(\pi)=3\) . Since the representation \(\wedge^{2}E\) identifies with the adjoint representation of \(\mathfrak{so}(q)\) (Lemma 6), it contains as a subrepresentation the complexification of the adjoint representation of G (Lemma 7). The condition \(\mathrm{M}_{4}(\pi)=3\) therefore implies that \(\mathfrak{g}_{(\mathbf{C})}=\mathfrak{so}(q)\) .

Conversely, suppose that \(\mathfrak{g}_{(\mathbf{C})} = \mathfrak{so}(q)\) . The adjoint representation of \(\mathfrak{so}(q)\) is irreducible, since \(\dim(\mathrm{E}) \neq 4\) (LIE, VIII, § 13, p. 193, (I) and p. 206, (I)), and of dimension \(n(n-1)/2\) . The representation \(\mathsf{S}_0^2(\mathrm{E})\) of \(\mathfrak{so}(q)\) has as highest weight \(2\varpi_1\) . By comparing its dimension with that of the irreducible representation of \(\mathfrak{so}(q)\) of highest weight \(2\varpi_1\) , computed by the Weyl formula (cf. LIE, VIII, §9, no. 2, th. 2 and LIE, VI, plates II and IV), one verifies that \(\mathsf{S}_0^2(\mathrm{E})\) is irreducible. We therefore conclude that \(\mathrm{M}_4(\pi) = 3\) if \(\mathfrak{g}_{(\mathbf{C})} = \mathfrak{so}(q)\) .

Remarks. --- 1) The condition “G is infinite” is necessary in Theorem 3 (exercise 9 of V, p. 507).

\begin{enumerate}
\def\labelenumi{\arabic{enumi})}
\setcounter{enumi}{1}
\item
  Let \(n \in \mathbf{N}\) . For every unitary representation \(\varrho\) of a compact group in a Hilbert space of dimension \(n\) , we have \(\mathrm{M}_{4}(\varrho) \leqslant n^{2}\) ; equality is possible (exercise 15 of V, p. 508).
\item
  One can prove (cf. R. GURALNICK and P.H. TIEP, Decomposition of small tensor powers and Larsen's conjecture, Representation Theory 9 (2005), 138--208) that the condition that G is infinite can be omitted if one assumes that E has dimension \(\geqslant 7\) and if one replaces the hypothesis \(M_{4}(\pi) = 2\) (resp. \(M_{4}(\pi) = 3\) ) by \(M_{8}(\pi) = 24\) (resp. \(M_{8}(\pi) = 105\) ).
\end{enumerate}

\specialchapter{Exercises}

\exercisesection{}

\begin{enumerate}
\def\labelenumi{\arabic{enumi})}
\item
  Give an example of a group G, a field K, and a faithful representation of G such that the associated \(K[G]\)-module is not faithful.
\item
  Let \(\pi\) be a unitary representation of G on a Hilbert space E and \(x \in E\) of norm 1 such that
\end{enumerate}

\[
\sup _ {g \in \mathrm{G}} \| \pi (g) x - x \| \leqslant 1.
\]

There exists a nonzero G-invariant vector \(y \in E\).

\begin{enumerate}
\def\labelenumi{\arabic{enumi})}
\setcounter{enumi}{2}
\tightlist
\item
  Let E be a Hilbert space. We denote by \(\mathbf{U}(\mathbf{E})\) the group of unitary endomorphisms of E and \(\mathbf{GL}(\mathbf{E})\) the group of invertible endomorphisms of E.
\end{enumerate}

\begin{enumerate}
\def\labelenumi{\alph{enumi})}
\item
  The group \(\mathbf{U}(\mathrm{E})\) endowed with the topology induced by the topology of simple convergence on \(\mathcal{L}(\mathrm{E})\) is a topological group, denoted \(\mathbf{U}(\mathrm{E})_{s}\) ; if E is infinite-dimensional, the group \(\mathbf{GL}(\mathrm{E})\) is not a topological group with the topology induced by the topology of simple convergence.
\item
  The topology of \(\mathbf{U}(\mathrm{E})_{s}\) coincides with the topology induced by the topology of compact convergence.
\item
  If E is countably generated, the group \(\mathbf{U}(\mathrm{E})_{s}\) is metrizable and complete.
\item
  Let L be the space of \(u \in \mathcal{L}(\mathrm{E})\) such that \(u^{*} = -u\) ; it is the Lie algebra of the real Lie group G = U(E) endowed with the topology induced by the Banach space topology of \(\mathcal{L}(\mathrm{E})\) (LIE, III, p. 146, cor. 2). The exponential map from L to \(\mathbf{U}(\mathrm{E})_s\) is continuous.
\item
  If E is of infinite dimension, then \(\mathbf{U}(\mathrm{E})_{s}\) is not a real Lie group of type \(\{L\}\) . (Consider the subgroup K of \(\mathbf{U}(\mathrm{E})_{s}\) formed by the unitary endomorphisms diagonal in a fixed orthonormal basis of E, and note that K is compact.)
\item
  The exponential map is not a local homeomorphism from L to \(\mathbf{U}(\mathbf{E})_{s}\) . (Consider its restriction to the diagonal endomorphisms in a fixed orthonormal basis.)
\end{enumerate}

\begin{enumerate}
\def\labelenumi{\arabic{enumi})}
\setcounter{enumi}{3}
\tightlist
\item
  Let G be a topological group. We say that G is amenable if there exists a linear form \(\lambda\colon\mathcal{C}_{b}(\mathrm{G};\mathbf{R})\to\mathbf{R}\) such that
\end{enumerate}

\begin{enumerate}
\def\labelenumi{(\roman{enumi})}
\item
  If \(f \geqslant 0\) , then \(\lambda(f) \geqslant 0\) .
\item
  One has \(\lambda(1)=1\) .
\item
  For every \(g \in \mathbf{G}\) and every \(f \in \mathcal{C}_b(\mathbf{G};\mathbf{R})\) , we have \(\lambda(\boldsymbol{\gamma}(g)f) = \lambda(\boldsymbol{\delta}(g)f)\) , where
\end{enumerate}

\[
\boldsymbol {\gamma} (g) f (x) = f \left(g ^ {- 1} x\right), \quad \boldsymbol {\delta} (g) f (x) = f (x g).
\]

If G is a discrete group, this definition coincides with that of EVT, IV, p. 73, exerc. 4.

\begin{enumerate}
\def\labelenumi{\alph{enumi})}
\tightlist
\item
  The topological group G is amenable if and only if for every integer \(n \in N\) , every finite family \((f_i)_{1 \leqslant i \leqslant n}\) in \(\mathcal{C}_b(\mathbf{G}; \mathbf{R})\) , all families \((g_i)_{1 \leqslant i \leqslant n}\) and \((h_i)_{1 \leqslant i \leqslant n}\) in G, every real number a, the condition
\end{enumerate}

\[
\sum_ {i = 1} ^ {n} (f _ {i} (x) - f _ {i} (g _ {i} x h _ {i})) \geqslant a
\]

for all \(x \in G\) implies that \(a \leqslant 0\) . (To show that this condition is sufficient, consider the set of functions of the form

\[
x \mapsto \sum_ {i = 1} ^ {n} (f _ {i} (x) - f _ {i} (g _ {i} x h _ {i}))
\]

and apply the Hahn--Banach theorem.)

\begin{enumerate}
\def\labelenumi{\alph{enumi})}
\setcounter{enumi}{1}
\item
  Every compact topological group is amenable, and every abelian topological group is amenable.
\item
  Let G be a non-commutative free group. The discrete group G is not amenable.
\item
  For every Hilbert space E and every continuous bounded representation of G in E, there exists a Hilbert norm q on E such that the representation of G in the Hilbert space E endowed with q is unitary.
\end{enumerate}

\begin{enumerate}
\def\labelenumi{\arabic{enumi})}
\setcounter{enumi}{4}
\tightlist
\item
  Let G be a topological group and H an open subgroup of G. We denote by p the canonical projection \(G \rightarrow G/H\) .
\end{enumerate}

\begin{enumerate}
\def\labelenumi{\alph{enumi})}
\item
  Let \(\sigma\colon\mathrm{G/H}\to\mathrm{G}\) be a map such that \(p\circ\sigma=\mathrm{Id}_{\mathrm{G/H}}\) . We denote by \(\beta\colon\mathrm{G}\times\mathrm{G/H}\to\mathrm{H}\) the map defined by \(\beta(g,x)=\sigma(gx)^{-1}g\) . The map \(\beta\) is continuous.
\item
  Let \(\pi\) be a bounded representation of H in a Hilbert space E. We denote by F the Hilbert space \(\ell^2 (\mathrm{G / H};\mathrm{E})\) . The formula
\end{enumerate}

\[
\varrho (g) f (x) = \pi (\beta (g ^ {- 1}, x) ^ {- 1}) f (g ^ {- 1} x)
\]

defines a continuous bounded representation of G in F.

\begin{enumerate}
\def\labelenumi{\alph{enumi})}
\setcounter{enumi}{2}
\tightlist
\item
  Suppose there exists a Hilbert norm q on F such that the representation \(\varrho\) is a unitary representation of G in \((F,q)\) . Then there exists a Hilbert norm \(\widetilde{q}\) on E such that \(\pi\) is a unitary representation of H in \((E,\widetilde{q})\) .
\end{enumerate}

\begin{enumerate}
\def\labelenumi{\arabic{enumi})}
\setcounter{enumi}{5}
\tightlist
\item
  \begin{enumerate}
  \def\labelenumii{\alph{enumii})}
  \tightlist
  \item
    Let E be a complex Banach space and let \(u \in \mathcal{L}(\mathrm{E})\) . If there exists \(n \in N\) such that \(n \geqslant 1\) and \(u^{n} = 1_{E}\) then the spectrum of u is contained in the group of n-th roots of unity and consists of eigenvalues of u. b) Let G be a commutative topological group such that every element of G has finite order. Every continuous irreducible representation of G in a Banach space has dimension 1, and takes values in the group \(U \subset C^{*}\) of complex numbers of modulus 1.
  \end{enumerate}
\end{enumerate}

\begin{enumerate}
\def\labelenumi{\alph{enumi})}
\setcounter{enumi}{2}
\tightlist
\item
  Let G be a commutative topological group such that the set of elements of finite order of G is dense in G. Every continuous irreducible representation of G in a Banach space has dimension 1, and takes values in the group \(U \subset C^{*}\) of complex numbers of modulus 1.
\end{enumerate}

\begin{enumerate}
\def\labelenumi{\arabic{enumi})}
\setcounter{enumi}{6}
\tightlist
\item
  For every integer \(n \geqslant 1\) , we denote by \(s_{n}\) the function \(Z \rightarrow R\) defined by
\end{enumerate}

\[
s _ {n} (k) = \left\{ \begin{array}{l l} (1 - k) ^ {n (1 - k)} & \text{if } k \leqslant - 1 \\ 1 & \text{if } k = 0 \\ (k + 1) ^ {- (k + 1) / n} & \text{if } k \geqslant 1. \end{array} \right.
\]

We denote by S the set of functions \(s_n\) .

\begin{enumerate}
\def\labelenumi{\alph{enumi})}
\item
  For every integer \(n \geqslant 1\) and every \((k_{1}, k_{2}) \in \mathbf{Z}^{2}\) , we have \(s_{4n}(k_{1})s_{4n}(k_{2}) \geqslant s_{n}(k_{1} + k_{2})\) .
\item
  For \(f \in \mathbf{C}(\mathrm{T})\) , we denote by \((a_k(f))_{k \in \mathbf{Z}}\) the coefficients of the Laurent expansion of \(f\) at 0. For every \(s \in \mathrm{S}\) and every \(f \in \mathbf{C}(\mathrm{T})\) , the series
\end{enumerate}

\[
\sum_ {k \in \mathbf {Z}} | a _ {k} (f) | s (k)
\]

converges, and the sets

\[
\mathcal {V} (s, \varepsilon) = \{f \in \mathbf {C} (\mathrm{T}) | \sum_ {k \in \mathbf {Z}} | a _ {k} (f) | s (k) <   \varepsilon \}
\]

defined for \(s \in S\) and \(\varepsilon > 0\) , form a fundamental system of open neighborhoods of 0 for a locally convex topology on \(\mathbf{C}(T)\) . The multiplication \(\mathbf{C}(T) \times \mathbf{C}(T) \to \mathbf{C}(T)\) is continuous for this topology.

\begin{enumerate}
\def\labelenumi{\alph{enumi})}
\setcounter{enumi}{2}
\tightlist
\item
  The topology defined in the preceding question is metrizable but not complete.
\end{enumerate}

\(d)\) For \(f \in \mathbf{C}(\mathrm{T})^*\) , we denote by \(\varrho(f)\) the map \(g \mapsto fg\) from \(\mathbf{C}(\mathrm{T})\) into \(\mathbf{C}(\mathrm{T})\) . The map \(\varrho\) is a continuous representation of \(\mathbf{C}(\mathrm{T})^*\) in the locally convex space \(\mathbf{C}(\mathrm{T})\) .

\begin{enumerate}
\def\labelenumi{\alph{enumi})}
\setcounter{enumi}{4}
\tightlist
\item
  The representation \(\varrho\) is irreducible.
\end{enumerate}

\exercisesection{}

In the exercises of this section, unless otherwise stated, G denotes a locally compact topological group, equipped with a left Haar measure which is denoted by \(\mu\) .

\begin{enumerate}
\def\labelenumi{\arabic{enumi})}
\item
  Let \(\pi\) be a unitary representation of G on a complex Hilbert space E. The representation \(\pi\) is irreducible if and only if the image of \(\mathcal{H}(\mathrm{G})\) in \(\mathcal{L}(\mathrm{E})\) is dense for the topology of pointwise convergence. (Use Exercise 24 of IV, p. 326.)
\item
  Let H be a closed subgroup of G. Let \(\nu\) be a quasi-invariant measure on G/H (INT, VII, p. 56, § 2, n° 5, th. 2), and let \(\varrho\) be the function on G with values in \(\mathbf{R}_{+}^{*}\) associated with it. The maps \(\varrho(g)\) defined by
\end{enumerate}

\[
\varrho (g) f (x) = \left(\frac {\varrho (g ^ {- 1} x)}{\varrho (x)}\right) ^ {1 / 2} f (g ^ {- 1} x)
\]

define a continuous linear representation of G on \(\mathcal{L}^{2}(\mathrm{G / H},\nu)\) , and, by passing to the quotient, define a unitary representation of G on \(\mathrm{L}^2 (\mathrm{G / H},\nu)\) , called the quasi-regular representation of \(\mathrm{G / H}\) .

\begin{enumerate}
\def\labelenumi{\arabic{enumi})}
\setcounter{enumi}{2}
\tightlist
\item
  Let G = R and let H be the group R equipped with the discrete topology and with the counting measure, which is a Haar measure.
\end{enumerate}

\begin{enumerate}
\def\labelenumi{\alph{enumi})}
\item
  For every \(t \in G\) , and every \(f \in L^{2}(H)\) , the map \(\varrho(t)f: x \mapsto f(x + t)\) belongs to \(L^{2}(H)\) ; the map \(\varrho(t)\) is a unitary map on \(L^{2}(H)\) . b) For all \(f,g \in L^{2}(H)\) , the map from G to C defined by \(t \mapsto \langle f | \varrho(t)g \rangle\) is measurable.
\item
  Let \(f\) and \(g\) be \(\mathrm{L}^2(\mathrm{H})\) . If there exists \(t \in \mathbf{R}\) such that \(\langle f | \varrho(t)g \rangle\) is nonzero, then the map \(t \mapsto \langle f | \varrho(t)g \rangle\) from \(\mathbf{G}\) into \(\mathbf{C}\) is not continuous. (In particular, the hypothesis that E is countably generated cannot be omitted from cor. 2 of INT, VIII, p. 171, §4, n° 7.)
\end{enumerate}

\(d)\) There exists \(f \in \mathrm{L}^2(\mathrm{H})\) such that the subspace of \(\mathrm{L}^2(\mathrm{H})\) generated by the elements \(\varrho(t)f\) , where \(t \in \mathbf{R}\) , is dense in \(\mathrm{L}^2(\mathrm{H})\) .

\begin{enumerate}
\def\labelenumi{\arabic{enumi})}
\setcounter{enumi}{3}
\item
  Let \(p \geqslant 1\) . The kernel of the biregular representation of G on \(\mathrm{L}^{p}(\mathrm{G}, \mu)\) is the image of the center of G under the homomorphism \(g \mapsto (g, g)\) .
\item
  Let G be a locally compact unimodular group.
\end{enumerate}

\begin{enumerate}
\def\labelenumi{\alph{enumi})}
\item
  If there exists a finite-dimensional square-integrable representation of G, then G is compact.
\item
  If there exists a square-integrable representation of G, then the center of G is compact.
\end{enumerate}

\begin{enumerate}
\def\labelenumi{\arabic{enumi})}
\setcounter{enumi}{5}
\tightlist
\item
  We assume that the group G is unimodular. We denote by \(\widehat{G}_{d} \subset \widehat{G}\) the set of classes of irreducible unitary square-integrable representations of G. For \(\pi \in \widehat{G}_{d}\) , we denote by \(p_{\pi}\) the orthogonal projector on \(L^{2}(G)\) whose image is the \(\pi\)-isotypic component of the right regular representation of G.
\end{enumerate}

\(a)\) Let \(f \in \mathcal{L}^2(\mathrm{G})\) and let \(\pi \in \widehat{\mathrm{G}}_d\) . We denote by E the representation space of \(\pi\) . The map \(g \mapsto f(g)\pi(g)\) belongs to \(\mathrm{L}^1(\mathrm{G};\mathcal{L}(\mathrm{E}))\) ; the endomorphism

\[
v = \int_ {\mathrm{G}} f (g) \pi (g) d \mu (g)
\]

is a Hilbert--Schmidt endomorphism, denoted \(v = \pi(f)\) .

\begin{enumerate}
\def\labelenumi{\alph{enumi})}
\setcounter{enumi}{1}
\item
  Let \(f \in \mathcal{L}^2(\mathrm{G})\) . We have \(\| \pi(\overline{f}) \|_2 = \| p_\pi(f) \|\) .
\item
  Let \(L_{d}^{2}(G)\) be the closed subspace of \(L^{2}(G)\) generated by the \(\pi\)-isotypic components of \(L^{2}(G)\) for \(\pi \in \widehat{G}_{d}\) . It is the Hilbert sum of the images of the orthogonal projectors \(p_{\pi}\) for \(\pi \in \widehat{G}_{d}\) . For every \(f \in L_{d}^{2}(G)\) , we have
\end{enumerate}

\[
\| f \| ^ {2} = \sum_ {\pi \in \widehat {\mathrm{G}} _ {d}} c (\pi) \| \pi (f) \| _ {2} ^ {2}
\]

where \(c(\pi)\) is the formal degree of \(\pi\) relative to \(\{e\}\) .

\begin{enumerate}
\def\labelenumi{\arabic{enumi})}
\setcounter{enumi}{6}
\tightlist
\item
  Let \(n \in N\) . Determine the unitary automorphic representations of \(R^{n}\) (that is, the representations \(\Gamma\) -automorphic for every discrete subgroup \(\Gamma\) of \(R^{n}\) ).
\end{enumerate}

¶ 8) We assume that G is a locally compact unimodular group. Let Γ be a discrete subgroup of G such that the quotient space X = Γ\textbackslash G is compact. Let μ be a Haar measure on G and β the counting measure on Γ; we denote by \(\widetilde{\mu} = \mu/\beta\) .

Let \(\varrho\) be the unitary representation of G on \(\mathrm{L}^{2}(\mathrm{X},\widetilde{\mu})\) defined in Example 3 of V, p. 404.

For \(\gamma\in\Gamma\) , we denote by \(\Gamma_{\gamma}\) (resp. \(\mathbf{G}_{\gamma}\) ) the centralizer of \(\gamma\) in \(\Gamma\) (resp. \(\mathbf{G}\) ). \(a)\) The group \(\mathbf{G}_{\gamma}\) is unimodular.

\begin{enumerate}
\def\labelenumi{\alph{enumi})}
\setcounter{enumi}{1}
\item
  Let \(\varphi\in\mathcal{K}(\mathrm{G})\) . For every \(\Gamma\)-automorphic unitary representation \(\pi\) , the endomorphism \(\pi(\varphi)\) has finite trace.
\item
  One has
\end{enumerate}

\[
\operatorname{Tr} (\varrho (\varphi)) = \sum_ {\pi \in \widehat {\mathrm{G}}} m _ {\pi} (\varrho) \operatorname{Tr} (\pi (\varphi))
\]

where \(m_{\pi}(\varrho)\) is the multiplicity of \(\pi\) in \(\varrho\) (“spectral side of the trace formula”).

\begin{enumerate}
\def\labelenumi{\alph{enumi})}
\setcounter{enumi}{3}
\tightlist
\item
  Let \(\Gamma^{\sharp}\) be the set of conjugacy classes of \(\Gamma\) . For \(\gamma \in \Gamma^{\sharp}\) , we denote by \(\beta_{\gamma}\) the counting measure on \(\Gamma_{\gamma}\) . We have
\end{enumerate}

\[
\mathrm{Tr} (\varrho (\varphi)) = \sum_ {\gamma \in \Gamma^ {\sharp}} \int_ {\Gamma_ {\gamma} \backslash G} \varphi (x ^ {- 1} \gamma x) d (\mu / \beta_ {\gamma}) (x).
\]

\begin{enumerate}
\def\labelenumi{\alph{enumi})}
\setcounter{enumi}{4}
\tightlist
\item
  We fix a Haar measure \(\mu_{\gamma}\) on \(G_{\gamma}\) for all \(\gamma \in \Gamma\) . We have
\end{enumerate}

\[
\operatorname{Tr} (\varrho (\varphi)) = \sum_ {\gamma \in \Gamma^ {\sharp}} \int_ {G _ {\gamma} \backslash G} \int_ {\Gamma_ {\gamma} \backslash G _ {\gamma}} \varphi (x ^ {- 1} y ^ {- 1} \gamma y x) d (\mu_ {\gamma} / \beta_ {\gamma}) (y) d (\mu / \mu_ {\gamma}) (x).
\]

\(f)\) For every \(\gamma \in \Gamma\) , the quotient measure \(\mu_{\gamma} / \beta_{\gamma}\) of \(\Gamma_{\gamma} \backslash G_{\gamma}\) is finite. \(g)\) One has

\[
\mathrm{Tr} (\varrho (\varphi)) = \sum_ {\gamma \in \Gamma^ {\sharp}} (\mu_ {\gamma} / \beta_ {\gamma}) (\Gamma_ {\gamma} \backslash \mathrm{G} _ {\gamma}) \int_ {\mathrm{G} _ {\gamma} \backslash \mathrm{G}} \varphi (x ^ {- 1} \gamma x) d (\mu / \mu_ {\gamma}) (x)
\]

(“geometric side of the trace formula”; the formula expressing the equality of this expression with that obtained in question c) is called the “trace formula” for Γ-automorphic representations).

\begin{enumerate}
\def\labelenumi{\alph{enumi})}
\setcounter{enumi}{7}
\tightlist
\item
  Set G = R and \(\Gamma = Z\) . Verify that the trace formula is equivalent in this case to the Poisson formula (cf. II, p. 239).
\end{enumerate}

\begin{enumerate}
\def\labelenumi{\arabic{enumi})}
\setcounter{enumi}{8}
\item
  We assume that G is a perfect real semisimple Lie group without compact factor. There is no nontrivial finite-dimensional unitary representation of G.
\item
  \begin{enumerate}
  \def\labelenumii{\alph{enumii})}
  \tightlist
  \item
    Every matrix coefficient of the left regular representation of G on \(L^{2}(G)\) belongs to \(\mathcal{C}_{0}(G)\) .
  \end{enumerate}
\end{enumerate}

\begin{enumerate}
\def\labelenumi{\alph{enumi})}
\setcounter{enumi}{1}
\tightlist
\item
  If G is not compact, the left regular representation of G on L²(G) contains no nonzero finite-dimensional subrepresentation.
\end{enumerate}

\begin{enumerate}
\def\labelenumi{\arabic{enumi})}
\setcounter{enumi}{10}
\tightlist
\item
  Let E be a Hilbert space.
\end{enumerate}

\begin{enumerate}
\def\labelenumi{\alph{enumi})}
\item
  Let \(\varrho\) be a unitary representation of the group \(\mathbf{R}\) in E. Let D be a dense subspace of E stable under \(\varrho\) such that, for every \(x\in\mathrm{D}\) , the map \(\psi_{x}:t\mapsto\varrho(t)x\) is differentiable on R. Let u be the partial operator with domain D such that \(u(x) = i^{-1}\psi_x'(0)\) for all \(x \in D\) . Then u is essentially self-adjoint and \(\overline{u}\) is the infinitesimal generator of \(\varrho\) .
\item
  Let u be a self-adjoint partial operator on E and let D be a dense subspace of \(\text{dom}(u)\) . Suppose that for every \(t \in R\) , we have \(\exp(itu)(D) \subset D\) . Then the space D is dense in \(E_{u}\) .
\item
  Let u be a self-adjoint partial operator on E. Let F be a Hilbert space and v a self-adjoint partial operator on F. The partial operator \(u \otimes 1_{F} + 1_{E} \otimes E\) on \(E \otimes F\) defines an essentially self-adjoint partial operator on \(E \widehat{\otimes}_{2} F\) .
\end{enumerate}

\begin{enumerate}
\def\labelenumi{\arabic{enumi})}
\setcounter{enumi}{11}
\tightlist
\item
  Let G be a real Lie group and let \(\varrho\) be a unitary representation of G in a Hilbert space E. We say that an element \(x \in \mathrm{E}\) is a vector of class \(\mathrm{C}^{\infty}\) if the map from G to E defined by \(g \mapsto \varrho(g)x\) is of class \(\mathrm{C}^{\infty}\) . a) Let \(\pi\) be a unitary representation of G in F and \(u: \mathrm{E} \to \mathrm{F}\) a G-morphism. For every vector \(x \in \mathrm{E}\) of class \(\mathrm{C}^{\infty}\) , the vector \(u(x)\) is of class \(\mathrm{C}^{\infty}\) in F.
\end{enumerate}

We denote by \(\mathrm{C}^{\infty}(\varrho)\) the space of vectors of class \(C^{\infty}\) in E. By question a), every G-morphism from \(\varrho\) into \(\pi\) defines, by passing to subspaces, a linear map from \(\mathrm{C}^{\infty}(\varrho)\) into \(\mathrm{C}^{\infty}(\pi)\) .

\begin{enumerate}
\def\labelenumi{\alph{enumi})}
\setcounter{enumi}{1}
\tightlist
\item
  The space \(\mathrm{C}^{\infty}(\varrho)\) is dense in E.
\end{enumerate}

\begin{enumerate}
\def\labelenumi{\arabic{enumi})}
\setcounter{enumi}{12}
\tightlist
\item
  Let \(\pi\) be an irreducible unitary representation of G in a complex Hilbert space E. Let \(\varrho\) be a unitary representation of G in a complex Hilbert space F. Suppose there exists a closed densely defined partial operator \(u\) from E into F whose domain is a G-stable subspace of E and which satisfies \(u \circ \pi(g) = \varrho(g) \circ u\) for all \(g \in \mathbf{G}\) .
\end{enumerate}

The partial operator \(u\) belongs to \(\mathcal{L}(\mathrm{E};\mathrm{F})\) , and there exists \(\lambda > 0\) such that \(\lambda u\) is isometric.

\begin{enumerate}
\def\labelenumi{\arabic{enumi})}
\setcounter{enumi}{13}
\tightlist
\item
  Let Z be a subgroup of the center C of G such that C/Z is compact and let \(\nu\) be a Haar measure on G/Z. Let \(\pi\) be an irreducible unitary representation of G in a complex Hilbert space E. We assume that \(\pi\) is square-integrable modulo the center. Let \(u \in \mathcal{L}(\mathrm{E})\) be a nuclear endomorphism (def. 4 of IV, p. 169).
\end{enumerate}

\begin{enumerate}
\def\labelenumi{\alph{enumi})}
\tightlist
\item
  For every \(x\) and \(y\) in E, the function
\end{enumerate}

\[
g \mapsto \langle x | \pi (g) u \pi (g ^ {- 1}) y \rangle
\]

belongs to \(\mathcal{F}(\mathrm{G / Z})\) .

\begin{enumerate}
\def\labelenumi{\alph{enumi})}
\setcounter{enumi}{1}
\tightlist
\item
  For every \(x\) and \(y\) in E, we have
\end{enumerate}

\[
\int_ {\mathrm{G} / \mathrm{Z}} \langle x \mid \pi (g) u \pi (g ^ {- 1}) y \rangle d \nu (g) = \frac {1}{c _ {\mathrm{Z}} (\pi)} \langle x \mid y \rangle \operatorname{Tr} (u).
\]

\begin{enumerate}
\def\labelenumi{\arabic{enumi})}
\setcounter{enumi}{14}
\tightlist
\item
  Let Z be a subgroup of the center C of G such that C/Z is compact and let \(\nu\) be a Haar measure on G/Z. Let \(\pi\) be an irreducible unitary representation of G in a complex Hilbert space E. Denote by \(\chi\) the restriction to Z of the central character of \(\pi\) .
\end{enumerate}

We denote by \(f_{x,y}\) the matrix coefficient \(g \mapsto \langle x | \pi(g)y \rangle\) for \((x, y) \in \mathrm{E}^{2}\) . Let \(p \geqslant 1\) be a real number and let \(q\) be the conjugate exponent of p.

\begin{enumerate}
\def\labelenumi{\alph{enumi})}
\item
  Let \(y \in \mathrm{E}\) such that the matrix coefficients \(f_{x,y}\) belong to \(\mathcal{L}_{\chi}^{p}(\mathrm{G})\) . The map \(x \mapsto f_{x,y}\) is a continuous G-morphism from \(\pi\) into \(\boldsymbol{\delta}_{\mathrm{G},\chi}^{(p)}\) .
\item
  Let \(\varrho\) be an irreducible unitary representation of G in a complex Hilbert space F whose central character restricts to \(\chi\) on \(Z\) . One assumes that \(p > 1\) and that the matrix coefficients of \(\pi\) belong to \(\mathcal{L}_{\chi}^{p}(\mathrm{G})\) and those of \(\varrho\) belong to \(\mathcal{L}_{\chi}^{q}(\mathrm{G})\) . Let \(x_0 \in \mathrm{E}\) and \(y_0 \in \mathrm{F}\) . For every \(x \in \mathrm{E}\) there exists a continuous linear form \(\ell_x\) from F to C such that
\end{enumerate}

\[
\ell_ {x} (y) = \int_ {\mathrm{G/Z}} \overline {{\langle x _ {0} | \pi (g) x \rangle}} \langle y _ {0} | \varrho (g) y \rangle d \nu (g)
\]

for all \(y \in \mathbf{F}\) .

\begin{enumerate}
\def\labelenumi{\alph{enumi})}
\setcounter{enumi}{2}
\item
  There exists a continuous linear map u from E into F such that \(\ell_{x}(y)=\langle u(x)\mid y\rangle\) for all \((x,y)\in\mathrm{E}\times\mathrm{F}\) .
\item
  One has
\end{enumerate}

\[
\int_ {\mathrm{G} / \mathrm{Z}} \overline {{\langle x | \pi (g) x ^ {\prime} \rangle}} \langle y | \varrho (g) y ^ {\prime} \rangle d \nu (g) = 0
\]

for \((x,x^{\prime},y,y^{\prime})\in \mathrm{E}^{2}\times \mathrm{F}^{2}\)

\begin{enumerate}
\def\labelenumi{\alph{enumi})}
\setcounter{enumi}{4}
\tightlist
\item
  We assume that p = 1. The matrix coefficients of \(\pi\) are orthogonal to those of \(\varrho\) .
\end{enumerate}

\begin{enumerate}
\def\labelenumi{\arabic{enumi})}
\setcounter{enumi}{15}
\tightlist
\item
  Let \(\mathbf{G} = \mathbf{S}\mathbf{L}(2, \mathbf{R})\) and let H be the open subset of C formed by the complex numbers z such that \(\mathcal{I}(z) > 0\) . Let \(\mu\) be the Lebesgue measure on C. For every element
\end{enumerate}

\[
g = \left( \begin{array}{c c} a & b \\ c & d \end{array} \right)
\]

of G and every \(z \in \mathbf{H}\) we set \(j(g, z) = cz + d\) .

\begin{enumerate}
\def\labelenumi{\alph{enumi})}
\tightlist
\item
  The group G acts continuously on the left on H by
\end{enumerate}

\[
\left( \begin{array}{c c} a & b \\ c & d \end{array} \right) \cdot z = \frac {a z + b}{c z + d}.
\]

The measure \(\nu = y^{-2} \cdot \mu\) is G-invariant.

\begin{enumerate}
\def\labelenumi{\alph{enumi})}
\setcounter{enumi}{1}
\tightlist
\item
  Let n be an integer such that \(n \geqslant 2\) . We denote by \(E_{n}\) the vector space of holomorphic functions \(f: H \to C\) such that the function \(z \mapsto |f(z)|^{2} \mathcal{I}(z)^{n}\) belongs to \(\mathcal{L}^{2}(\mathbf{H}, \nu)\) . The space \(E_{n}\) endowed with the sesquilinear form
\end{enumerate}

\[
\langle f \mid g \rangle = \int_ {\mathbf {H}} \overline {{f (z)}} g (z) \mathcal {I} (z) ^ {n} d \nu (z)
\]

is a Hilbert space of countable type and of infinite dimension.

\begin{enumerate}
\def\labelenumi{\alph{enumi})}
\setcounter{enumi}{2}
\tightlist
\item
  For every \(g \in G\) and every \(f \in E_{n}\) we set
\end{enumerate}

\[
\pi_ {n} (g) f (z) = j (g ^ {- 1}, z) ^ {- n} f (g ^ {- 1} \cdot z).
\]

The map \(\pi_n\) defines a unitary representation of G in \(\mathrm{E}_n\) . \(d)\) The function \(f_n\) defined by \(f_n(z) = (z + i)^{-n}\) belongs to \(\mathrm{E}_n\) ; it satisfies \(\pi_n(k)f_n = \chi_{-n}(k)f_n\) for all \(k \in \mathbf{SO}(\mathbf{R}^2)\) , where \(\chi_n\) is the character such that

\[
\chi_ {n} \Big (\left( \begin{array}{c c} \cos (\theta) & \sin (\theta) \\ - \sin (\theta) & \cos (\theta) \end{array} \right) \Big) = e ^ {i n \theta}
\]

for all \(\theta\in\mathbf{R}\) .

\begin{enumerate}
\def\labelenumi{\alph{enumi})}
\setcounter{enumi}{4}
\item
  Let F be a subrepresentation of \(E_{n}\) . The orthogonal projector of \(E_{n}\) whose image is the \(\chi_{-n}\)-isotypic component of the restriction of \(\pi_{n}\) to \(\mathbf{SO}(\mathbf{R}^{2})\) is given by \(f \mapsto -(2i)^{n}f(i)f_{n}\) . (Use cor. 2 of V, p. 465 and the Cauchy integral formula, FVC, to appear.)
\item
  The representation \(\pi_n\) is irreducible.
\end{enumerate}

\(g)\) The matrix coefficient \(g \mapsto \langle f_n | \pi_n(g)f_n \rangle\) is nonzero and belongs to \(\mathcal{L}^2(\mathrm{G})\) .

\begin{enumerate}
\def\labelenumi{\alph{enumi})}
\setcounter{enumi}{7}
\tightlist
\item
  The representation \(\pi_{n}\) is square-integrable and there exists a constant c > 0 such that the formal degree of \(\pi_{n}\) is \(c(n - 1)\) .
\end{enumerate}

\begin{enumerate}
\def\labelenumi{\arabic{enumi})}
\setcounter{enumi}{16}
\item
  Let G be a connected real Lie group. There exists a faithful finite-dimensional unitary representation of G if and only if the group G is locally isomorphic to a compact group.
\item
  Let X be a locally compact space equipped with a bounded positive measure \(\mu\) . For every Banach space F, we denote by S(X, \(\mu\) ; F) the space of classes of functions \(\mu\) -measurable from X into F modulo the relation of equality \(\mu\) -almost everywhere, endowed with the topology of convergence in measure (INT, IV, p. 194, § 5, no. 11). We denote by A the commutative subalgebra of \(\mathcal{L}(\mathrm{L}_{\mathbf{C}}^{2}(\mathrm{X},\mu))\) formed by the multiplication endomorphisms by functions \(f \in \mathcal{L}_{\mathbf{C}}^{\infty}(\mathrm{X},\mu)\) .
\end{enumerate}

We denote by \(e\colon \mathbf{R}\to \mathbf{C}\) the homomorphism \(t\mapsto \exp (2i\pi t)\) .

Let E be a topological vector space.

\begin{enumerate}
\def\labelenumi{\alph{enumi})}
\tightlist
\item
  Let \(f \in \mathrm{S}(\mathrm{X}, \mu; \mathbf{R})\) such that \(|f(x)| \geqslant 1\) for \(\mu\) -almost all \(x \in \mathrm{X}\) . There exists a real number \(t \in ]0, 1[\) such that
\end{enumerate}

\[
\mu (\{x \in \mathrm{X} \mid t f (x) \in [ 1 / 4, 3 / 4 ] + \mathbf {Z} \}) \geqslant \frac {2}{5} \mu (\mathrm{X}).
\]

\begin{enumerate}
\def\labelenumi{\alph{enumi})}
\setcounter{enumi}{1}
\item
  Let u be a linear map from E into \(\mathrm{S}(\mathrm{X},\mu;\mathbf{R})\) . For every \(x\in E\) , we denote by \(\varrho_{u}(x)\) the element of A given by the endomorphism of multiplication by the function \(e\circ u(x)\) , which belongs to \(\mathcal{L}_{\mathbf{C}}^{\infty}(\mathrm{X},\mu)\) . The map \(\varrho_{u}\) is a continuous representation of G into \(\mathrm{L}_{\mathbf{C}}^{2}(\mathrm{X},\mu)\) if and only if the map u is continuous. (If u is not continuous, use the previous question to verify that \(\varrho\) is not continuous.)
\item
  Let \(\varrho\) be a unitary representation of E in \(\mathrm{L}^2 (\mathrm{X},\mu)\) whose image is contained in \(\mathcal{A}\) . There exists a unique continuous linear map \(u\colon \mathrm{E}\to \mathrm{S}(\mathrm{X},\mu ;\mathbf{R})\) such that \(\varrho = \varrho_{u}\) . (First consider the case where \(\mathrm{E} = \mathbf{R}\) , and then apply the Stone theorem.)
\end{enumerate}

\begin{enumerate}
\def\labelenumi{\arabic{enumi})}
\setcounter{enumi}{18}
\item
  Let G be a finite group and H a subgroup of G. For every unitary representation \(\pi\) of H in a finite-dimensional complex Hilbert space E, prove that the induced unitary representation \(\operatorname{Ind}_{\mathrm{H}}^{\mathrm{G}}(\pi)\) is isomorphic to the induced representation defined in A, VIII, p. 392, § 21, n° 3.
\item
  Let G be a locally compact group which is countable at infinity, and let H be a closed subgroup of G. Let \(\pi\) be a unitary representation of H in a Hilbert space E and let \(\varrho\) be the unitary representation of G induced by \(\pi\) . We denote by F the representation space of \(\varrho\) .
\end{enumerate}

We fix Haar measures \(\mu\) on G and \(\beta\) on H, and we denote by \(\kappa\colon\mathbf{G}\to\mathbf{R}_{+}^{*}\) a continuous function such that \(\kappa(xh)=\Delta_{\mathrm{H}}(h)\Delta_{\mathrm{G}}(h)^{-1}\kappa(x)\) for \((x,h)\in\mathbf{G}\times\mathrm{H}\) and such that \(\kappa(e)=1\) .

\begin{enumerate}
\def\labelenumi{\alph{enumi})}
\item
  If G/H admits a bounded G-invariant measure, and if there exists a nonzero vector \(x \in E\) which is H-invariant, then there exists a nonzero G-invariant vector \(y \in F\).
\item
  Conversely, suppose that \(y \in F\) is a nonzero and G-invariant vector of \(\varrho\) . Show that there exists a function \(f \in L_{\overline{\pi}}^{2}\) nonzero such that
\end{enumerate}

\[
f (x) = \left(\frac {\kappa (g ^ {- 1} x)}{\kappa (x)}\right) ^ {1 / 2} f (g ^ {- 1} x)
\]

for all \((g, x) \in \mathrm{G}^{2}\) . Deduce that \(\kappa = 1\) , and then the measure \(\mu/\beta\) on G/H is defined and bounded, and finally that \(f(e) \in \mathrm{E}\) is a nonzero H-invariant vector.